\section{Phase One --- Consumer Economy through GameFi}

\subsection{Why GameFi is the entry vector}

Three properties make GameFi the right Phase-One surface for an alliance
whose end-state is a financial super-app and a multi-vertical compounding
network.

\textbf{Regulatory geometry.} Gaming transactions are not securities,
banking, or money transmission per se. They sit in a regulatory regime
(consumer protection, anti-fraud, payment processing) that is materially
easier to operate under in launch than a direct financial product. The
Alliance can ship at scale on day one without queueing for financial
licenses.

\textbf{Transaction volume per user.} A live GameFi user generates
hundreds to thousands of on-chain interactions per month --- microtransactions,
NFT mints, marketplace trades, tournament entries, reward claims. This is
two to three orders of magnitude more transaction throughput per user than
a typical fintech or consumer financial product. High throughput
\emph{validates the substrate}, \emph{stresses the operations posture},
and \emph{generates protocol fees} from day one.

\textbf{User-base acquisition economics.} Gaming user-acquisition costs
are dominated by content, marketing, and live-ops rather than regulatory
overhead. The fifteen-title slate (\S 2.2) brings its own user base on
launch; the Alliance does not start from zero subscribers.

\subsection{Mechanics of Phase One}

The fifteen-title slate ships across mobile, web, and desktop targets.
Each title is instrumented to:

\begin{itemize}[leftmargin=2em,itemsep=0.1em]
\item Use the W3A L2 as its transaction substrate (high-frequency
microtransactions clear sub-second with finality on the L2; finality on
the parent chain follows the standard rollup cadence).
\item Issue and trade in-game assets as NFTs (ERC-721 / ERC-1155) and
fungible in-game currency (LP-curated ERC-20 with constrained
issuance schedule).
\item Settle player payments through the partner stack: card and
stablecoin payments via SF Private Pay on the US/CA side; UK / EU /
APAC payments via the regional partners; remittance corridors via SSB
and M-Pesa for the diaspora flows.
\item Provide a unified player wallet with optional self-custody
(WalletConnect) or custodial (W3A MPC + KMS) modes. Default is custodial
for the casual user; self-custody is a single-click opt-in for the
sovereignty-minded user.
\item Publish a single-sign-on identity (Hanzo IAM tenant per game studio
under the standard \texttt{<org>-<app>} convention) that carries
forward into Phase Two without re-onboarding.
\end{itemize}

\subsection{What success looks like at end of Phase One}

Phase-One success is measured against transaction throughput, monthly
active wallets, gross revenue from in-game economies, and qualified-user
conversion eligibility for Phase Two. Concrete targets:

\begin{itemize}[leftmargin=2em,itemsep=0.1em]
\item \textbf{500,000 monthly active wallets} across the fifteen titles
within twelve months of launch.
\item \textbf{One million on-chain transactions per day} (sustained
30-day average) by month nine.
\item \textbf{\$25 million annualised gross revenue} from in-game
economies and marketplace fees by end of month twelve.
\item \textbf{200,000 wallets KYC-eligible} for Phase-Two conversion at
end of month twelve (i.e.\ users who have passed a baseline accredited
or non-accredited KYC and are eligible to transact in the financial
super-app of \S 4 without an onboarding gap).
\end{itemize}

These are the metrics that justify the Phase-Two unlock. If the Alliance
clears each of them within the targeted window, the conversion to the
financial super-app of \S 4 proceeds; if any is missed, the Phase-Two
ramp is held until the gating metric is met. The discipline is to
\emph{not skip phases} in pursuit of headline metrics that don't carry
the user base forward.

\subsection{Substrate-validation as a parallel outcome}

A second outcome of Phase One, independent of the consumer-economy
metrics, is \emph{substrate validation at scale}. Running one million
on-chain transactions per day sustained against the W3A L2 produces an
operational and security record that no synthetic benchmark or testnet
campaign can match. That record is itself a Phase-Three asset: it is the
substrate proof point that institutional Phase-Three counterparties
diligence before signing.
